% chktex-file 1
% chktex-file 3
% chktex-file 8
% chktex-file 36
% Suppress chktex false positives:
%   1  -- \skewt uses \xspace so trailing space is intentional
%   3  -- math like \{Y_t\}_{t=1}^{n_t} parses correctly
%   8  -- en-dashes for proper-name pairs and ranges; ISO dates with hyphen
%   36 -- AR(p) is canonical notation; no space before the parenthesis
\documentclass[11pt]{article}

\usepackage[margin=1in]{geometry}
\usepackage{amsmath,amssymb,amsthm}
\usepackage{bm}
\usepackage{xspace}
\usepackage{booktabs}
\usepackage{enumitem}
\usepackage{listings}
\usepackage{xcolor}
\usepackage{hyperref}
\usepackage[capitalize]{cleveref}
\hypersetup{colorlinks=true, linkcolor=blue, urlcolor=blue, citecolor=blue}

% ---- Notation aligned with skewt_rev_2.tex / mycommands.sty ----
\newcommand{\bY}{\mathbf{Y}}
\newcommand{\bZero}{\boldsymbol{0}}
\newcommand{\skewt}{skew-$t$\xspace}
\newcommand{\Skewt}{Skew-$t$\xspace}
\newcommand{\eref}[1]{(\ref{#1})}
\newcommand{\sref}[1]{Section~\ref{#1}}

% ---- Local shorthands ----
\newcommand{\Lexact}{L_{\mathrm{exact}}}
\newcommand{\Lcond}{L_{\mathrm{cond}}}
\newcommand{\bphi}{\boldsymbol{\varphi}}

% ---- Theorem environments ----
\theoremstyle{plain}
\newtheorem{proposition}{Proposition}[section]
\newtheorem{theorem}[proposition]{Theorem}
\newtheorem{lemma}[proposition]{Lemma}
\newtheorem{corollary}[proposition]{Corollary}

\theoremstyle{definition}
\newtheorem{definition}[proposition]{Definition}

\theoremstyle{remark}
\newtheorem{remark}[proposition]{Remark}

\lstset{
    basicstyle=\ttfamily\small,
    keywordstyle=\color{blue},
    commentstyle=\color{gray}\itshape,
    stringstyle=\color{teal},
    showstringspaces=false,
    breaklines=true,
    frame=single,
    framesep=4pt,
    language=R
}

\title{The $\bphi$-Block of the AR(2) Sampler Targets the Wrong Full
       Conditional:\\ Exact versus Conditional Likelihood in
       \texttt{updatePhiAR2TS}}
\author{Yi-Xuan Xie}
\date{2026-07-13}

\begin{document}
\maketitle

\begin{abstract}
\noindent
The Metropolis--Hastings update for $\bphi = (\varphi_1, \varphi_2)$ in
\texttt{update\_params\_ar2.R} evaluates the Box--Jenkins \emph{conditional}
likelihood, while the latent-block updaters evaluate the \emph{exact}
stationary likelihood. The two Gibbs blocks therefore target different
distributions; this note proves the discrepancy, explains why the analogous
omission is exactly harmless in the original AR(1) code, quantifies the effect
at $n_t = 50$, and states the three-line fix.
\end{abstract}

\section{Setup and notation}\label{sec:setup}

\begin{itemize}
  \item Fix one latent channel ($z^\star$, $\tau^\star$, or a knot coordinate)
        and one knot $k$; write $Y_{1:T}$ for its copula-transformed series,
        with independent replicates across $k = 1, \dots, K$.
  \item The AR(2) prior on $Y_{1:T}$ is the stationary causal Gaussian AR(2)
        normalised to unit marginal variance, which by Yule--Walker pins
        \begin{equation}\label{eq:moments}
          \gamma_0 = 1, \qquad
          \gamma_1 = \frac{\varphi_1}{1 - \varphi_2}, \qquad
          \gamma_2 = \varphi_1 \gamma_1 + \varphi_2, \qquad
          \sigma_\epsilon^2 = 1 - \varphi_1 \gamma_1 - \varphi_2 \gamma_2,
        \end{equation}
        all functions of $\bphi = (\varphi_1, \varphi_2)$ alone
        (\texttt{get\_ar2\_standard\_moments}).
  \item Stationarity restricts $\bphi$ to the triangle
        \begin{equation}\label{eq:triangle}
          S_2 = \bigl\{ (\varphi_1, \varphi_2) :
            \varphi_1 + \varphi_2 < 1,\;
            \varphi_2 - \varphi_1 < 1,\;
            |\varphi_2| < 1 \bigr\}.
        \end{equation}
  \item Under stationarity the joint prior density factorises as
        \begin{equation}\label{eq:factor}
          p(Y_{1:T} \mid \bphi)
          = \underbrace{p(Y_1, Y_2 \mid \bphi)}_{\text{initial block}}
            \prod_{t=3}^{T}
            N\!\bigl( Y_t ;\, \varphi_1 Y_{t-1} + \varphi_2 Y_{t-2},\,
                      \sigma_\epsilon^2 \bigr),
        \end{equation}
        with initial block
        \begin{equation}\label{eq:init}
          (Y_1, Y_2)^\top \sim
          N_2\!\left( \bZero,\; \Sigma(\bphi) \right),
          \qquad
          \Sigma(\bphi) =
          \begin{pmatrix} 1 & \gamma_1 \\ \gamma_1 & 1 \end{pmatrix}.
        \end{equation}
  \item Equivalently, $p(Y_1) = N(0, 1)$ and
        \begin{equation}\label{eq:p2}
          p(Y_2 \mid Y_1, \bphi)
          = N\!\bigl( \gamma_1 Y_1,\; 1 - \gamma_1^2 \bigr),
        \end{equation}
        the ``Fix~1'' branch of \texttt{get\_ar2\_conditional\_params}
        (see \texttt{ar2\_fix\_proposal.tex}).
\end{itemize}

\section{Exact versus conditional likelihood}\label{sec:defs}

\begin{definition}[Exact and conditional likelihoods]\label{def:lik}
For fixed $Y_{1:T}$, define as functions of $\bphi \in S_2$
\begin{equation}\label{eq:defs}
  \Lexact(\bphi) = p(Y_{1:T} \mid \bphi),
  \qquad
  \Lcond(\bphi) = p(Y_{3:T} \mid Y_{1:2}, \bphi)
  = \prod_{t=3}^{T}
    N\!\bigl( Y_t ;\, \varphi_1 Y_{t-1} + \varphi_2 Y_{t-2},\,
              \sigma_\epsilon^2 \bigr).
\end{equation}
\end{definition}

\begin{itemize}
  \item By \eref{eq:factor} the two are linked through the discrepancy factor
        \begin{equation}\label{eq:discrepancy}
          \Lexact(\bphi) = f(\bphi;\, Y_{1:2}) \cdot \Lcond(\bphi),
          \qquad
          f(\bphi;\, Y_{1:2})
          = N_2\!\bigl( (Y_1, Y_2)^\top ;\, \bZero,\, \Sigma(\bphi) \bigr).
        \end{equation}
  \item \texttt{ar2\_transition\_loglik} computes $\log \Lcond$ only: it sums
        transition densities over \texttt{data[, 3:nt]}, discarding
        $\log f$ entirely.
\end{itemize}

\section{Why the omission is exact in AR(1)}\label{sec:ar1}

\begin{lemma}[AR(1) exemption]\label{lem:ar1}
For the unit-variance stationary AR(1), i.e.\ $p = 1$ with
$Y_1 \sim N(0, \gamma_0)$ and $\gamma_0 \equiv 1$, the discrepancy factor
satisfies $f(\varphi; Y_1) = N(Y_1; 0, 1)$, free of $\varphi$; hence
\begin{equation}\label{eq:ar1prop}
  \Lexact(\varphi) \propto_{\varphi} \Lcond(\varphi)
  \qquad \text{exactly, for every realisation } Y_{1:T}.
\end{equation}
\end{lemma}

\begin{proof}
The initial block for $p = 1$ is the single marginal
$p(Y_1 \mid \varphi) = N(Y_1; 0, \gamma_0)$, and the normalisation
$\gamma_0 = 1$ removes its only $\varphi$-dependence; \eref{eq:discrepancy}
with constant $f$ gives \eref{eq:ar1prop}.
\end{proof}

\begin{remark}[The original code is correct]\label{rem:orig}
The comment in \texttt{biometrics/update\_params.R} (line 925),
``\emph{likelihood of data impacted by phi does not include day 1},'' is
therefore an identity under the unit-variance normalisation, not an
approximation. \texttt{updatePhiTS} and the AR(1) latent blocks consequently
share one target posterior, and the original sampler is a valid
MH-within-Gibbs scheme.
\end{remark}

\section{The AR(2) violation}\label{sec:violation}

\begin{itemize}
  \item The sampler is MH-within-Gibbs on $(Y, \bphi)$ with target
        \begin{equation}\label{eq:target}
          \pi(Y, \bphi \mid \text{data})
          \;\propto\;
          p(\text{data} \mid Y)\; p(Y \mid \bphi)\; p(\bphi),
        \end{equation}
        and validity requires every block's acceptance ratio to be computed
        under this single $\pi$.
  \item The correct full conditional for $\bphi$ is, by
        \eref{eq:factor}--\eref{eq:init},
        \begin{equation}\label{eq:fullcond}
          \pi(\bphi \mid Y)
          \;\propto\;
          \underbrace{N_2\!\bigl( (Y_1, Y_2)^\top ;\, \bZero,\,
            \Sigma(\bphi) \bigr)}_{f(\bphi;\, Y_{1:2})}
          \prod_{t=3}^{T}
          N\!\bigl( Y_t ;\, \varphi_1 Y_{t-1} + \varphi_2 Y_{t-2},\,
                    \sigma_\epsilon^2 \bigr)
          \; p(\bphi).
        \end{equation}
\end{itemize}

\begin{proposition}[Target mismatch between the two blocks]\label{prop:mismatch}
Suppose $T \geq 3$ and $p(\bphi) > 0$ on $S_2$. In
\texttt{update\_params\_ar2.R}:
\begin{enumerate}[label=(\roman*)]
  \item the $Y$-block updaters (\texttt{updateZTS\_AR2},
        \texttt{updateTauTS\_AR2}, \texttt{updateKnotsTS\_AR2}) evaluate the
        prior contribution of $Y_2$ through \eref{eq:p2}, hence are
        $\pi$-invariant MH kernels;
  \item the $\bphi$-block (\texttt{updatePhiAR2TS}) accepts with ratio built
        from $\Lcond(\bphi)\, p(\bphi)$, hence is invariant for
        \begin{equation}\label{eq:tildepi}
          \tilde{\pi}(Y, \bphi \mid \text{data})
          \;\propto\;
          p(\text{data} \mid Y)\;
          N(Y_1; 0, 1)\, \Lcond(\bphi)\; p(\bphi),
        \end{equation}
        and $\tilde{\pi} \neq \pi$ because
        $\pi / \tilde{\pi} \propto f(\bphi; Y_{1:2}) / N(Y_2; 0, 1)$ is
        non-constant in $\bphi$ whenever $\gamma_1 \neq 0$;
  \item the composite kernel therefore does not admit $\pi$ as an invariant
        distribution.
\end{enumerate}
\end{proposition}

\begin{proof}
Claim (i) is read off the ``Fix~1'' $t = 2$ branch used in each latent
update.

For (ii), $\gamma_1 = \varphi_1 / (1 - \varphi_2)$ makes $\Sigma(\bphi)$,
hence $f(\bphi; Y_{1:2})$, non-constant on $S_2$ for any $Y_{1:2}$ with
$Y_1 Y_2 \neq 0$ or $Y_1^2 + Y_2^2 \neq 2$: explicitly,
\begin{equation}\label{eq:logf}
  \log f(\bphi;\, Y_{1:2})
  = -\log 2\pi
    - \tfrac{1}{2} \log\bigl( 1 - \gamma_1^2 \bigr)
    - \frac{Y_1^2 - 2 \gamma_1 Y_1 Y_2 + Y_2^2}
           {2 \bigl( 1 - \gamma_1^2 \bigr)}
\end{equation}
depends on $\bphi$ through $\gamma_1$.

The factor $f$ depends jointly on $\bphi$ and on $Y_{1:2}$, so it can be
absorbed into the normalising constant of neither block's full conditional.

For (iii), suppose the composite kernel had invariant law $\pi$; since the
blocks touch disjoint coordinates, each factor kernel would then be
$\pi$-invariant on its own coordinates. But the $\bphi$-kernel's unique
invariant conditional is
$\tilde{\pi}(\bphi \mid Y) \neq \pi(\bphi \mid Y)$, a contradiction.
\end{proof}

\begin{remark}[Provenance]\label{rem:port}
The port copied the ``skip the head of the series'' shape of the AR(1) code
(\cref{rem:orig}) without noticing that the skipped term acquires
$\bphi$-dependence at $p = 2$. Only the $Y_1$ factor stays exempt: it is
$N(0, 1)$ by the $\gamma_0 = 1$ normalisation for every $p$.
\end{remark}

\section{Magnitude at the study's sample size}\label{sec:magnitude}

\begin{itemize}
  \item The omitted factor \eref{eq:logf} carries the information of $2$ of
        the $T$ observations per series, so the perturbation of the
        $\bphi$-posterior is $O(T^{-1})$ with constant governed by
        $\partial \log f / \partial \bphi$.
  \item A direct Monte Carlo check ($T = 50$ as in the simulation study,
        $30$ replicate datasets, componentwise MH with the production
        $N(0, 0.5^2)$ prior and silent rejection on $S_2$) compares the
        posterior mean under $\Lcond$ against $\Lexact$;
        \cref{tab:bias} reports the averages.
\end{itemize}

\begin{table}[ht]
\centering
\caption{Posterior means of $(\varphi_1, \varphi_2)$ under the conditional
         and exact likelihoods; $T = 50$, $30$ replicates, $K$ independent
         series per dataset. Settings are the fixed-AR(2) DGPs of
         \texttt{simstudy/setup.R}.}
\label{tab:bias}
\begin{tabular}{llccc}
\toprule
Setting & $K$ & Truth & $\Lcond$ (code) & $\Lexact$ (with $f$) \\
\midrule
16 near-unit-root & 1 & $(0.15,\; 0.80)$  & $(0.16,\; 0.76)$  & $(0.15,\; 0.77)$ \\
16 near-unit-root & 5 & $(0.15,\; 0.80)$  & $(0.14,\; 0.81)$  & $(0.13,\; 0.82)$ \\
9 strong          & 1 & $(0.80,\, -0.35)$ & $(0.73,\, -0.27)$ & $(0.73,\, -0.27)$ \\
9 strong          & 5 & $(0.80,\, -0.35)$ & $(0.78,\, -0.33)$ & $(0.78,\, -0.33)$ \\
10 weak           & 1 & $(0.12,\, -0.05)$ & $(0.06,\, -0.10)$ & $(0.06,\, -0.10)$ \\
10 weak           & 5 & $(0.12,\, -0.05)$ & $(0.11,\, -0.04)$ & $(0.12,\, -0.04)$ \\
\bottomrule
\end{tabular}
\end{table}

\begin{remark}[Interpretation]\label{rem:magnitude}
All $\Lcond$-versus-$\Lexact$ differences are at most $0.02$, within Monte
Carlo noise, while the visible departures from truth ($0.03$--$0.07$) are the
ordinary $O(T^{-1})$ finite-sample AR bias present under both likelihoods.
The sampler is thus theoretically invalid but numerically indistinguishable
from the corrected one at $T = 50$; the gap scales as $2/T$ and becomes
worth re-checking if $T$ drops to $10$--$20$ (time-block forecasting).
\end{remark}

\section{The fix}\label{sec:fix}

\begin{itemize}
  \item Add the $\bphi$-dependent part of $\log f$, namely
        \eref{eq:p2} evaluated at the data, to both \texttt{cur.ll} and
        \texttt{can.ll} in \texttt{updatePhiAR2TS} (each with its own
        moments):
\end{itemize}

\begin{lstlisting}
# inside updatePhiAR2TS, after computing (cur|can).moments:
g1 <- moments$gamma1  # = phi1 / (1 - phi2)
ll <- ll + sum(dnorm(data.t2, g1 * data.t1,
                     sqrt(1 - g1^2), log = TRUE))
# data.t1, data.t2 = the t = 1, 2 slices of `data`
# (columns for day.mar = 2; third-margin slices for day.mar = 3)
\end{lstlisting}

\begin{itemize}
  \item The $Y_1$ marginal $N(0, 1)$ needs no term: it cancels in the
        acceptance ratio for every $p$ (\cref{rem:port}).
  \item The silent-rejection proposal scheme is unaffected: Proposition~3.3
        of \texttt{where\_we\_are.tex} concerns the proposal's invariance on
        $S_2$ and holds for whichever target density is supplied.
  \item Expected effect: none on the $T = 50$ simulation conclusions
        (\cref{tab:bias}); the patch restores \eref{eq:fullcond} as the
        $\bphi$-block target and hence validity of the composite sampler.
\end{itemize}

\section{Verification and status}\label{sec:verify}

\begin{itemize}
  \item \textbf{Status (2026-07-15): applied.} The patch of \sref{sec:fix}
        is in production as \texttt{ar2\_stationary\_loglik}
        (\texttt{auxfunctions\_ar2.R}), called at all three likelihood
        evaluations inside \texttt{updatePhiAR2TS}
        (\texttt{update\_params\_ar2.R}).
  \item \textbf{Scope of the correctness claim.} The proof of
        \cref{prop:mismatch} and its resolution never reference the
        data-generating process, so validity of the fixed sampler holds for
        every DGP unconditionally; only the \emph{negligibility} claims below
        are DGP- and $T$-specific.
\end{itemize}

\subsection{Geweke prior-recovery test (DGP-free)}\label{sec:geweke}

\begin{itemize}
  \item Design (successive-conditional simulator, Geweke 2004): switch the
        data likelihood off, so the target collapses to
        $p(\bphi)\, p(Y \mid \bphi)$, and alternate the production
        $Y$-block prior structure ($t$, $t{+}1$, $t{+}2$ companion terms via
        \texttt{ar2\_conditional\_log\_density}) with the production
        \texttt{updatePhiAR2TS}. If and only if the blocks share one joint
        density must the $\bphi$-marginal of the chain equal the
        $N(0, 0.5^2)^{\otimes 2}$ prior truncated to $S_2$ (reference:
        rejection sample, $n \approx 10^6$).
  \item Small $T$ amplifies the discrepancy: the initial block carries
        $2/T$ of the information, so $T = 6$ and $T = 4$ were used
        ($1/3$ and $1/2$ respectively).
\end{itemize}

\begin{table}[ht]
\centering
\caption{Geweke prior-recovery: deviation of the chain's $\bphi$-marginal
         mean from the truncated-prior reference, in MCMC standard errors
         ($z$). ``Conditional'' = pre-fix production code; ``stationary'' =
         patched code.}
\label{tab:geweke}
\begin{tabular}{llcc}
\toprule
Design & Version & $z\bigl(E[\varphi_1]\bigr)$ & $z\bigl(E[\varphi_2]\bigr)$ \\
\midrule
$T=6$, $K=3$, one chain of $115$k
  & conditional & $1.0$ & $\phantom{-}2.1$ \\
  & stationary  & $0.4$ & $-0.7$ \\
\addlinespace
$T=4$, $K=3$, six chains of $40$k
  & conditional & $0.5$ & $\phantom{-}1.8$ \\
  & stationary  & $0.1$ & $-0.6$ \\
\bottomrule
\end{tabular}
\end{table}

\begin{itemize}
  \item The conditional-only version deviates on $E[\varphi_2]$ in the
        \emph{same direction} in two independent designs (deviations
        $+0.021$ and $+0.011$; Stouffer-combined $z \approx 2.8$), matching
        the theory: $f$ acts through
        $\gamma_1 = \varphi_1 / (1 - \varphi_2)$ and so distorts the
        $\varphi_2$-geometry of the truncated prior first.
  \item The stationary version passes every statistic in both designs
        ($|z| \leq 1$), confirming the patch removes the incompatibility in
        the production function itself, not merely in a reimplementation.
\end{itemize}

\subsection{DGP coverage of the negligibility claim}\label{sec:coverage}

\begin{itemize}
  \item \cref{tab:bias} covers the fixed-AR(2) DGPs (settings 9--16);
        \cref{tab:coverage} extends the same comparison ($T = 50$, $K = 1$,
        $30$ replicates) to the remaining DGP families of
        \texttt{simstudy/setup.R}.
  \item For the long-memory rows the reference is the Yule--Walker
        pseudo-true projection \texttt{arfima\_yw\_projection(d)\$ar2}; the
        residual gap to it is finite-sample bias, identical under both
        likelihoods.
\end{itemize}

\begin{table}[ht]
\centering
\caption{Conditional-versus-exact posterior means on the remaining DGP
         families; last column is the largest coordinatewise difference.}
\label{tab:coverage}
\begin{tabular}{lcccc}
\toprule
DGP & Pseudo-true & $\Lcond$ & $\Lexact$ & $\max$ diff \\
\midrule
White noise (settings 1--8) & $(0.00,\, 0.00)$ & $(0.02,\, 0.01)$ & $(0.02,\, 0.01)$ & $0.004$ \\
ARFIMA $d = 0.20$ (17)      & $(0.22,\, 0.11)$ & $(0.18,\, 0.08)$ & $(0.18,\, 0.08)$ & $0.001$ \\
ARFIMA $d = 0.35$ (18)      & $(0.42,\, 0.21)$ & $(0.34,\, 0.19)$ & $(0.34,\, 0.19)$ & $0.001$ \\
ARFIMA $d = 0.45$ (19)      & $(0.58,\, 0.29)$ & $(0.59,\, 0.25)$ & $(0.58,\, 0.25)$ & $0.003$ \\
\bottomrule
\end{tabular}
\end{table}

\subsection{Permanent regression guard}\label{sec:guard}

\begin{itemize}
  \item \texttt{tests/test\_phi\_block\_consistency.R} now protects the
        property in three layers: a structural assertion
        (\texttt{updatePhiAR2TS} calls \texttt{ar2\_stationary\_loglik} and
        never \texttt{ar2\_transition\_loglik}), an exact identity check
        (the added term equals the latent-block $t = 2$ density
        \eref{eq:p2} to $10^{-10}$ on $200$ random draws), and a fixed-seed
        Geweke smoke gate ($|z| < 4$).
  \item Layer two is the theorem-level guard --- it verifies block
        compatibility directly and costs milliseconds; the heavier
        detection of \cref{tab:geweke} is documented here rather than
        re-run per test invocation.
\end{itemize}

\end{document}
